\note{sa page 2}

sur $\mathcal{E}^{(\nu)}$ une structure d'\struck{algèbre} anneau gradué ?

Il suffit que l'on ait les relations (1.1.)

\[ \begin{cases} \pi_{i}^{2} = \pi_{i} \\ \pi_{i} \pi_{j} = \pi_{j} \pi_{i} = 0 \ \text{ si } i \neq j \\ \sum_{i} \pi_{i} = 1 \end{cases} \]

qui impliquent que les $\mathrm{pr}_{\nu}^{\mathcal{E}}$ définis par (1.2.) sont
des projecteurs deux à deux orthogonaux de somme $1$, donc que les
$\mathcal{E}^{(\nu)}$ définissent bien une graduation \add{de $k$-modules} sur
$\mathcal{E}^{*}$, \struck{puis} et que
$\mathcal{E}^{(\nu)} \mathcal{E}^{(\nu')} \subset \mathcal{E}^{(\nu + \nu')}$.

\note{Le renvoi « (1.2.) » est écrit ainsi ; la formule qui définit
$\mathrm{pr}_{\nu}$ porte, page 2, l'étiquette (1.3) b). La discordance est sur
le feuillet.}

Noter aussi que (1.4) sont valables, \struck{en \ill{}} avec \ill{}
\struck{où $u = \sum_{i,j} \pi_{j} u \pi_{i}$} \struck{de la relation
$\pi_{i+\nu} u \pi_{i} = \pi_{i} u \pi_{i}$}

\[ \begin{cases} \mathcal{E} = \bigoplus_{i,j} \pi_{j}\, \mathcal{E}\, \pi_{i} \\ \mathrm{pr}_{i,j}(u) = \pi_{j}\, u\, \pi_{i} \end{cases} \qquad \text{où } u = \sum_{\alpha,\beta} \pi_{\beta}\, u\, \pi_{\alpha} \]

et la relation $u \pi_{i} = \pi_{i+\nu} u$ implique, puisque
$u \pi_{i} = \sum_{\beta} \pi_{\beta} u \pi_{i}$ et
$\pi_{i+\nu} u = \sum_{\alpha} \pi_{i+\nu} u \pi_{\alpha}$, que
\[ \pi_{\beta}\, u\, \pi_{i} = 0 \quad \text{si } \beta \neq i + \nu . \]

\marginal{Regarder les modules sur \struck{des} l'anneau $\mathcal{E}$ muni des éléments $\pi_{i}$ comme donnés.}

3. \struck{Nous pouvons} Partons enfin d'un \struck{algèbre} anneau gradué
$\mathcal{E}^{*}$ \struck{sur $k$}, et posons-nous la question de l'existence et
de l'unicité d'une famille $(\pi_{i})_{i \in \mathbb{Z}}$ \add{satisfaisant
(1.1.)} d'\ill{} de $\mathcal{E}^{*}$, définissant la graduation, par (1.3.)
(ou par (1.4)). On suppose ici

